\chapter{Przykłady elementów pracy dyplomowej}

\section{Liczba}

Pakiet \texttt{siunitx} zadba o to, by liczba została poprawnie sformatowana: \\
\begin{center}
	\num{1234567890.0987654321}
\end{center}


\section{Rysunek}

Pakiet \texttt{subcaption} pozwala na umieszczanie w podpisie rysunku odnośników do ,,podilustracji'': \\

\begin{figure}[h]
	\centering
	\begin{subfigure}{0.35\textwidth}
		\centering
		\framebox[2.0\width]{A}
		\subcaption{\label{subfigure_a}}
	\end{subfigure}
	\begin{subfigure}{0.35\textwidth}
		\centering
		\framebox[2.0\width]{B}
		\subcaption{\label{subfigure_b}}
	\end{subfigure}
	
	\caption{\label{fig:subcaption_example} Przykład użycia \texttt{\textbackslash subcaption}: \protect\subref{subfigure_a} litera A, \protect\subref{subfigure_b} litera B.}
\end{figure}

\begin{figure}
\centering

\begin{tabular}{cccc}
\subfloat[Object detection]{\includegraphics[width = 1.5in]{imgs/obj_bbox.jpg}} &
\subfloat[Object segmentation]{\includegraphics[width = 1.5in]{imgs/obj_mask.jpg}} \\
\subfloat[Drivable area segmentation]{\includegraphics[width = 1.5in]{imgs/drivable.jpg}} &
\subfloat[Lane segmentation]{\includegraphics[width = 1.5in]{imgs/lanes.jpg}} 
\end{tabular}
% \caption{Przykłady zbiorów uczących. Zbiór (b) wygenerowano tylko dla sieci MultiTask V3, a zbiory (a),(c) i (d) dla wszystkich modeli.}
\caption{Examples of training sets. Set (b) was generated for the MultiTask V3 network only, and sets (a), (c) and (d) for all models.}
\label{fig:example_dataset}
\end{figure}


\begin{figure}[!]	
\begin{adjustwidth}{0cm}{0cm}
% \centering %% If there is a figure in wide page, please release command \centering
\includegraphics[width=1\linewidth]{imgs/fast_example.pdf}
\end{adjustwidth}
\caption{Determining the feature point flag for an exemplary image patch.
In that case none of the 16 pixels on the Bresenham circle have an intensity above $I_p + t$ ($229$ for $t=20$), while 14 of them have an intensity below $I_p - t$ ($189$).
This results in two 16-bit vectors, which are compared with 16 different patterns, each consisting of 16 bits: 9 ones and 7 zeros.
The vector representing pixels with intensity higher than $I_p + t$ does not match any pattern, while the second one matches 6 different patterns in total. Therefore, we determine the centre pixel is a~feature point (it fulfils the FAST's ``lower than'' condition). }
\label{fig:fast_example_flag}
\end{figure}



\section{Tabela}

Pakiet \texttt{threeparttable} umożliwia dodanie do tabeli adnotacji: \\

\begin{table}[h]
	\centering
	
	\begin{threeparttable}
		\caption{Przykład tabeli}
		\label{tab:table_example}
		
		\begin{tabularx}{0.6\textwidth}{C{1}}
			\toprule
			\thead{Nagłówek\tnote{a}} \\
			\midrule
			Tekst 1 \\
			Tekst 2 \\
			\bottomrule
		\end{tabularx}
		
		\begin{tablenotes}
			\footnotesize
			\item[a] Jakiś komentarz\textellipsis
		\end{tablenotes}
		
	\end{threeparttable}
\end{table}


\begin{table}[!t]
\setlength{\tabcolsep}{6pt}
\centering
\caption{Comparison of resource consumption on the Kria KV260 platform.}
\begin{tabular}{|c|c|c|c|c|c|c|c|}
\hline
\multirow{2}{*}{Resource usage} &  \multicolumn{4}{c|}{CPU cores} & \multirow{2}{*}{RAM} & \multirow{2}{*}{CMA} & \multirow{2}{*}{Power}\\
\cline{2-5}
{} & CPU$_0$ & CPU$_1$ & CPU$_2$ & CPU$_3$ & {} & {} & {}\\ 
\hline
Used & 85 \%  & 22 \%  & 3 \% & 3 \%  & 38 \%  & 6 \%  & 4.95 W\\
\hline
\end{tabular}
\label{tab:consumption}
\end{table}


\section{Wzory matematyczne}

Czasem zachodzi potrzeba wytłumaczenia znaczenia symboli użytych w równaniu. Można to zrobić z użyciem zdefiniowanego na potrzeby niniejszej klasy środowiska \texttt{eqwhere}.

\begin{equation}
E = mc^2
\end{equation}
gdzie
\begin{eqwhere}[2cm]
	\item[$m$] masa
	\item[$c$] prędkość światła w próżni
\end{eqwhere}

Odległość półpauzy od lewego marginesu należy dobrać pod kątem najdłuższego symbolu (bądź listy symboli) poprzez odpowiednie ustawienie parametru tego środowiska (domyślnie: 2 cm).


\begin{equation}
\label{eq:moments_mod_orient_approx}
	\begin{aligned}
		\tan\theta_{i} & \leq & \tan\theta (x,y) & \leq \tan\theta_{i+1} \\
		m_{10}(x,y) \tan\theta_{i} & \leq & m_{01} (x,y) & \leq m_{10} (x,y) \tan\theta_{i+1}
	\end{aligned}
\end{equation}
where:
$\theta(x,y)$ is the feature orientation, 
$m_{10}(x,y), m_{01}(x,y)$ are the image moments, 
$\tan\theta_{i}$, $\tan\theta_{i+1}$ are the tangent values of two adjacent intervals to which the orientation will be assigned, 10-bit fixed-point. 


\begin{equation}
\label{eq:BRIEF test}
 \tau (\textbf{p}; u_i, v_i ) := 
 \begin{cases} 
 1 & \text{for $I(u_i) < I(v_i)$} \\
 0 & \text{for $I(u_i) \ge I(v_i)$}
 \end{cases}
\end{equation}
where 
\textbf{p} is the smoothed image patch around a~given feature point $(u_i, v_i)$, 
$u_i, v_i$ are the coordinates of sample pairs \textbf{p} $(u_i \neq v_i)$, 
$I(w)$ is the brightness of the image in a~point $w = (x, y)$.


\begin{equation}
\label{eq:sum_corner_score}
V = \max \left ( \sum_{x \in S_{bright}} \left | I_{p \to x} - I_p \right | - t, \sum_{x \in S_{dark}} \left | I_{p} - I_{p \to x} \right | - t \right )
\end{equation}
\begin{equation}
\begin{aligned}
S_{bright} &= \left \{ x|I_{p \to x} \ge I_p + t \right \} \\
S_{dark} &= \left \{ x|I_{p \to x} \le I_p - t \right \}
\end{aligned}
\end{equation}

\begin{equation}\label{eq:hog_mod_amplituda_SRA}
	\begin{aligned}
		m(x,y) & = \sqrt{G_{x}(x,y)^{2} + G_{y} (x,y)^{2}} \\ & \approx \max ((0.875a + 0.5b), a) \\
		a      & = \max(\left| G_{x} (x,y)\right| ,\left| G_{y} (x,y)\right| )               \\
		b      & = \min(\left| G_{x} (x,y)\right| ,\left| G_{y} (x,y)\right| )
	\end{aligned}
\end{equation}